\FloatBarrier
\subsection{Measuring adaptation maturity}
\label{app:public-stock-adaptation}
Large search gains may partly repair errors that further policy learning
would remove. We therefore test longer adaptation before interpreting a
search comparison as robustness to stronger baseline policies. The complete
changed-opponent score criterion permits this separate study; it does not
remove the failed work-admissibility condition.

All six retained systems receive exactly 384 additional PPO updates from
their 128-update weights, with fresh Adam and the original teacher frozen.
This is a weight warm start, not restoration of optimizer or environment
state. Each update uses 512 environments and 64 decisions, with the same
architecture, terminal-match objective, teacher KL penalty and other PPO
settings as the earlier adaptation. Only the final checkpoint is eligible;
intermediate checkpoints cannot be selected or training extended.

Additional work totals 75,497,472 decisions, 3,536,901 rounds, 485,230 matches
and 2,304 updates across six systems. Each system has 16,777,216 cumulative
adaptation decisions; original parent training remains separate. Frozen-teacher
identity, initial/final checkpoint lineage, finite parameters and counters
are audited. Training counts are recorded, not independently replayed, and
training completion is not strength evidence.

Qualification uses 1,024 new duplicate streams, both seats, with zero online
search. Each new actor faces both retained old parents under its rule; the
old actor also faces the other old parent on the same streams. The own-parent
contrast is new score minus 0.5; deterministic duplicate self-play gives the
old actor exactly 0.5. The other-parent contrast subtracts the old actor's
score against that same opponent, preserving within-stream covariance.
All twelve contrasts share 20,000 bootstrap replicates and one-sided lower
quantile $0.05/12$. A rule qualifies only when both contrasts have positive
lower bounds for both parents. Control and at least one finalist must qualify
to launch the prospective all-improved-policy search study.

\begin{table}[ht]
\centering\small
\begin{tabular}{llrrr}
\toprule
Rule & Parent & New/own (\%) & New/other (\%) & Old/other (\%) \\
\midrule
Control & 0 & 49.10 & 50.12 & 49.17 \\
Control & 1 & 50.59 & 49.95 & 50.83 \\
Preview + choice & 0 & 57.52 & 57.28 & 48.85 \\
Preview + choice & 1 & 57.03 & 57.03 & 51.15 \\
All three & 0 & 57.71 & 59.47 & 50.15 \\
All three & 1 & 56.54 & 55.10 & 49.85 \\
\bottomrule
\end{tabular}

\caption{All zero-search qualification scores, with ties worth one half.
``Own'' and ``other'' always identify retained old opponents. The old actor
against itself has duplicate score 50\%; the old/other column is measured
on the same streams as new/other.}
\label{tab:public-stock-adaptation-scores}
\end{table}

The eighteen endpoints contain 36,864 matches, 264,776 rounds and 6,024,229
scored decisions, all replayed in Python. Independent reduction verifies
outcomes and counts; the recorded search counters are zero. Both finalists
qualify. Against their own retained actors, preview plus choice improves by
7.52/7.03 points and all three by 7.71/6.54 points. Their other-opponent
improvements are 8.42/5.88 and 9.33/5.25 points, respectively. All eight
associated lower bounds are positive.

Control does not qualify: its own-parent changes are $-0.90/+0.59$ points
and other-parent changes $+0.95/-0.88$, with all four lower bounds negative.
The all-improved-policy search launch condition therefore fails, and that
prospective study is not run. This is a failure to establish stronger control
policies, not evidence that search has lost its value. The two finalists now
have measured stronger actors; qualification measures no new-policy search
gain. A separately specified follow-up retains old control and tests the
qualified finalist actors, with unequal adaptation budgets and new cohorts
explicit.
The failed launch condition is neither relaxed nor relabeled as a pass.
\begin{table}[ht]
\centering\small
\begin{tabular}{llrr}
\toprule
Rule & Parent & Own change [LB] & Other change [LB] \\
\midrule
Control & 0 & -0.90 [-2.76] & +0.95 [-1.39] \\
Control & 1 & +0.59 [-1.32] & -0.88 [-3.17] \\
Preview + choice & 0 & +7.52 [+5.00] & +8.42 [+4.98] \\
Preview + choice & 1 & +7.03 [+4.54] & +5.88 [+2.47] \\
All three & 0 & +7.71 [+5.05] & +9.33 [+5.86] \\
All three & 1 & +6.54 [+3.83] & +5.25 [+1.73] \\
\bottomrule
\end{tabular}

\caption{All twelve paired qualification contrasts, in percentage points,
with one-sided lower bounds at quantile $0.05/12$. The nominal 95\%
familywise target is within this study only. A longer training duration
does not itself qualify an improved policy.}
\label{tab:public-stock-adaptation-contrasts}
\end{table}
\FloatBarrier
